% sec2_7.tex — Section 2.7: Testing Conditional Homoskedasticity

With the advent of robust standard errors allowing us to do inference without specifying
the conditional second moment $\E(\varepsilon_i^2 \mid \mathbf{x}_i)$, testing conditional
homoskedasticity is not as important as it used to be.  This section presents only the
most popular test due to White (1980) for the case of random samples.\footnote{See,
e.g., Judge et al.\ (1985, Section~11.3) or Greene (1997, Section~12.3) for other tests.}

Recall that $\widehat{\mathbf{S}}$ given in \eqref{eq:2.5.1} (with
$\hat{\varepsilon}_i = e_i$) is consistent for $\mathbf{S}$ (Proposition~2.4), and
$s^2 \mathbf{S}_{\mathbf{xx}}$ is consistent for $\sigma^2 \bm{\Sigma}_{\mathbf{xx}}$
(an implication of Proposition~2.2).  But under conditional homoskedasticity,
$\mathbf{S} = \sigma^2 \bm{\Sigma}_{\mathbf{xx}}$ (see \eqref{eq:2.6.4}), so the
difference between the two consistent estimators should vanish:
\begin{equation}
  \widehat{\mathbf{S}} - s^2 \mathbf{S}_{\mathbf{xx}}
    = \frac{1}{n} \sum_{i=1}^{n} e_i^2\, \mathbf{x}_i \mathbf{x}_i'
    - s^2 \frac{1}{n} \sum_{i=1}^{n} \mathbf{x}_i \mathbf{x}_i'
    = \frac{1}{n} \sum_{i=1}^{n} (e_i^2 - s^2)\, \mathbf{x}_i \mathbf{x}_i'
    \pto \mathbf{0}.
  \label{eq:2.7.1}
\end{equation}
Let $\bm{\psi}_i$ be a vector collecting unique and nonconstant elements of the
$K \times K$ symmetric matrix $\mathbf{x}_i \mathbf{x}_i'$.  (Construction of
$\bm{\psi}_i$ from $\mathbf{x}_i \mathbf{x}_i'$ will be illustrated in Example~2.7
below.)  Then \eqref{eq:2.7.1} implies
\begin{equation}
  \mathbf{c}_n \equiv \frac{1}{n} \sum_{i=1}^{n}
    (e_i^2 - s^2)\, \bm{\psi}_i \pto \mathbf{0}.
  \label{eq:2.7.2}
\end{equation}
This $\mathbf{c}_n$ is a sample mean converging to zero.  Under some conditions
appropriate for a Central Limit Theorem to be applicable, we would expect
$\sqrt{n}\,\mathbf{c}_n$ to converge in probability to a normal distribution with mean
zero and some asymptotic variance $\mathbf{B}$, so for any consistent estimator
$\widehat{\mathbf{B}}$ of $\mathbf{B}$,
\begin{equation}
  n \cdot \mathbf{c}_n' \widehat{\mathbf{B}}^{-1} \mathbf{c}_n
    \dto \chi^2(m),
  \label{eq:2.7.3}
\end{equation}
where $m$ is the dimension of $\mathbf{c}_n$.  For a certain choice of
$\widehat{\mathbf{B}}$, this statistic can be computed as $nR^2$ from the following
auxiliary regression:
\begin{equation}
  \text{regress } e_i^2 \text{ on a constant and } \bm{\psi}_i.
  \label{eq:2.7.4}
\end{equation}

White (1980) rigorously developed this argument\footnote{The original statement of
White's theorem is more general as it covers the case where $\{y_i, \mathbf{x}_i\}$ is
independent but not identically distributed (i.n.i.d.).} to prove

\begin{proposition}[White's Test for Conditional Heteroskedasticity]
\label{prop:2.6}
\emph{In addition to Assumptions~2.1 and 2.4, suppose that
\textup{(a)} $\{y_i, \mathbf{x}_i\}$ is i.i.d.\ with finite
$\E(\varepsilon_i^2 \mathbf{x}_i \mathbf{x}_i')$
\textup{(}thus strengthening Assumptions~2.2 and 2.5\textup{)},
\textup{(b)} $\varepsilon_i$ is independent of $\mathbf{x}_i$
\textup{(}thus strengthening Assumption~2.3 and conditional homoskedasticity\textup{)},
and \textup{(c)} a certain condition holds on the moments of $\varepsilon_i$ and
$\mathbf{x}_i$.  Then,}
\[
  nR^2 \dto \chi^2(m),
\]
\emph{where $R^2$ is the $R^2$ from the auxiliary regression \eqref{eq:2.7.4}, and $m$
is the dimension of $\bm{\psi}_i$.}
\end{proposition}

\begin{example}[regressors in White's $nR^2$ test]
\label{ex:2.7}
Consider the Cobb--Douglas cost function of Section~1.7:
\[
  \log\!\Bigl(\frac{TC_i}{p_{i3}}\Bigr)
    = \beta_1 + \beta_2 \log(Q_i) + \beta_3 \log\!\Bigl(\frac{p_{i1}}{p_{i3}}\Bigr)
    + \beta_4 \log\!\Bigl(\frac{p_{i2}}{p_{i3}}\Bigr) + \varepsilon_i.
\]
Here, $\mathbf{x}_i' = (1, \log(Q_i), \log(p_{i1}/p_{i3}), \log(p_{i2}/p_{i3}))$, a
four-dimensional vector.  There are $10$ $(= 4 \cdot 5/2)$ unique elements in
$\mathbf{x}_i \mathbf{x}_i'$:
\begin{gather*}
  1,\; \log(Q_i),\; \log\!\Bigl(\frac{p_{i1}}{p_{i3}}\Bigr),\;
    \log\!\Bigl(\frac{p_{i2}}{p_{i3}}\Bigr), \\[4pt]
  [\log(Q_i)]^2,\;
    \log(Q_i) \cdot \log\!\Bigl(\frac{p_{i1}}{p_{i3}}\Bigr),\;
    \log(Q_i) \cdot \log\!\Bigl(\frac{p_{i2}}{p_{i3}}\Bigr), \\[4pt]
  \biggl[\log\!\Bigl(\frac{p_{i1}}{p_{i3}}\Bigr)\biggr]^{\!2},\;
    \log\!\Bigl(\frac{p_{i1}}{p_{i3}}\Bigr)
      \cdot \log\!\Bigl(\frac{p_{i2}}{p_{i3}}\Bigr),\;
  \biggl[\log\!\Bigl(\frac{p_{i2}}{p_{i3}}\Bigr)\biggr]^{\!2}.
\end{gather*}
$\bm{\psi}_i$ is a nine-dimensional vector excluding a constant from this list.
So the $m$ in Proposition~2.6 is~9.
\end{example}

If White's test accepts the null of conditional homoskedasticity, then the results of
Section~2.6 apply, and statistical inference can be based on the $t$- and $F$-ratios from
Chapter~1.  Otherwise inference should be based on the robust $t$ and Wald statistics of
Proposition~2.3.

Because the regressors $\bm{\psi}_i$ in the auxiliary regression have many elements
consisting of squares and cross-products of the elements of $\mathbf{x}_i$, the test
should be consistent (i.e., the power approaches unity as $n \to \infty$) against most
heteroskedastic alternatives but may require a fairly large sample to have power close to
unity.  If the researcher knows that some of the elements of $\bm{\psi}_i$ do not affect
the conditional second moment, then they can be dropped from the auxiliary regression and
the power might be increased in finite samples.  The downside of it, of course, is that
if such knowledge is false, the test will have no power against heteroskedastic
alternatives that relate the conditional second moment to those elements that are excluded
from the auxiliary regression.
